\subsubsection*{Solution}
\paragraph*{Step-by-Step Derivation}

Define

\[
z=\frac{4\alpha}{(1+\alpha)^2}, \qquad H(n,z)={}_2F_1\left(\frac{1-n}{2},1-\frac n2;2;z\right),
\]

so that

\[
f(n,\alpha)=(1+\alpha)^{n-1}H(n,z).
\]

We first treat $0<\alpha<1$, for which $0<z<1$.

\subparagraph*{1. Euler integral representation}

Euler's integral representation of the Gaussian hypergeometric function is given in \href{\detokenize{https://dlmf.nist.gov/15.6.E1}}{DLMF Section 15.6.1}. Here it gives

\[
H(n,z) = \frac{1}{ \Gamma\left(1-\frac n2\right) \Gamma\left(1+\frac n2\right)} \int_0^1 t^{-n/2}(1-t)^{n/2}(1-zt)^{-(1-n)/2}\,dt.
\]

Let

\[
P(n)= \frac{1}{ \Gamma\left(1-\frac n2\right) \Gamma\left(1+\frac n2\right)}.
\]

Using $\psi(x)=\Gamma'(x)/\Gamma(x)$,

\[
\frac{P'(n)}{P(n)} = \frac12\psi\left(1-\frac n2\right) -\frac12\psi\left(1+\frac n2\right).
\]

Consequently,

\[
P(0)=1, \qquad P'(0)=0.
\]

For the integrand,

\[
\Phi(n,t)=t^{-n/2}(1-t)^{n/2}(1-zt)^{-(1-n)/2},
\]

and hence

\[
\frac{\partial}{\partial n}\log\Phi(n,t) = -\frac12\log t +\frac12\log(1-t) +\frac12\log(1-zt).
\]

At $n=0$,

\[
\Phi(0,t)=\frac1{\sqrt{1-zt}}.
\]

Therefore,

\[
H_n'(0,z) = \frac12\int_0^1 \frac{ \log\left(\dfrac{(1-t)(1-zt)}{t}\right) }{ \sqrt{1-zt} }\,dt. \tag{1}
\]

Differentiation under the integral is justified because, in a neighborhood of $n=0$, the integrand and its $n$-derivative are bounded by integrable functions proportional to

\[
t^{-\varepsilon}(1-t)^{-\varepsilon} \left(1+|\log t|+|\log(1-t)|\right),
\]

while $1-zt\geq 1-z>0$.

\subparagraph*{2. Evaluate $H(0,z)$}

At $n=0$,

\[
H(0,z) = \int_0^1\frac{dt}{\sqrt{1-zt}} = \frac{2(1-\sqrt{1-z})}{z}.
\]

Now

\[
1-z = 1-\frac{4\alpha}{(1+\alpha)^2} = \frac{(1-\alpha)^2}{(1+\alpha)^2}.
\]

Because $0\leq\alpha\leq1$,

\[
\sqrt{1-z}=\frac{1-\alpha}{1+\alpha}.
\]

Thus, for $\alpha>0$,

\[
H(0,z) = \frac{ 2\left(1-\dfrac{1-\alpha}{1+\alpha}\right) }{ \dfrac{4\alpha}{(1+\alpha)^2} } =1+\alpha. \tag{2}
\]

It follows that $f(0,\alpha)=1$.

\subparagraph*{3. Evaluate the logarithmic integral}

Introduce

\[
q=\sqrt{1-z}=\frac{1-\alpha}{1+\alpha}, \qquad z=1-q^2.
\]

In (1), set

\[
s=\sqrt{1-zt}, \qquad t=\frac{1-s^2}{z}, \qquad dt=-\frac{2s}{z}\,ds.
\]

The endpoints become

\[
t=0\Longrightarrow s=1, \qquad t=1\Longrightarrow s=q.
\]

Moreover,

\[
1-t=\frac{s^2-q^2}{z}, \qquad 1-zt=s^2.
\]

Therefore,

\[
\frac{(1-t)(1-zt)}{t} = \frac{s^2(s^2-q^2)}{1-s^2}.
\]

Equation (1) becomes

\[
H_n'(0,z) = \frac1z \int_q^1 \log\left( \frac{s^2(s^2-q^2)}{1-s^2} \right)ds. \tag{3}
\]

Write the integral in (3) as

\[
A(q)= 2\int_q^1\log s\,ds +\int_q^1\log(s-q)\,ds +\int_q^1\log(s+q)\,ds -\int_q^1\log(1-s)\,ds -\int_q^1\log(1+s)\,ds.
\]

Using

\[
\int \log x\,dx=x\log x-x,
\]

the individual terms are

\[
2\int_q^1\log s\,ds = -2-2q\log q+2q,
\]

\[
\int_q^1\log(s-q)\,ds = (1-q)\log(1-q)-(1-q),
\]

\[
\int_q^1\log(s+q)\,ds = (1+q)\log(1+q)-(1+q)-2q\log(2q)+2q,
\]

\[
-\int_q^1\log(1-s)\,ds = -(1-q)\log(1-q)+(1-q),
\]

and

\[
-\int_q^1\log(1+s)\,ds = -2\log2+2+(1+q)\log(1+q)-(1+q).
\]

Adding and simplifying,

\[
A(q) = 2\left[ q-1 +(1+q)\log\left(\frac{1+q}{2}\right) -2q\log q \right]. \tag{4}
\]

Since

\[
q-1=-\frac{2\alpha}{1+\alpha}, \qquad 1+q=\frac{2}{1+\alpha}, \qquad \frac{1+q}{2}=\frac1{1+\alpha},
\]

equation (4) becomes

\[
A(q) = -\frac4{1+\alpha} \left[ \alpha+\log(1+\alpha) +(1-\alpha)\log\left(\frac{1-\alpha}{1+\alpha}\right) \right]. \tag{5}
\]

Because $H_n'(0,z)=A(q)/z$ and

\[
z=\frac{4\alpha}{(1+\alpha)^2},
\]

we obtain

\[
\frac{H_n'(0,z)}{1+\alpha} = -\frac1\alpha \left[ \alpha+\log(1+\alpha) +(1-\alpha)\log\left(\frac{1-\alpha}{1+\alpha}\right) \right]. \tag{6}
\]

\subparagraph*{4. Differentiate $f(n,\alpha)$}

The product rule gives

\[
g(\alpha) = \frac1{1+\alpha} \left[ \log(1+\alpha)H(0,z)+H_n'(0,z) \right].
\]

Using $H(0,z)=1+\alpha$ from (2),

\[
g(\alpha) = \log(1+\alpha)+\frac{H_n'(0,z)}{1+\alpha}.
\]

Substituting (6),

\[
g(\alpha) = \log(1+\alpha) -\frac1\alpha \left[ \alpha+\log(1+\alpha) +(1-\alpha)\log\left(\frac{1-\alpha}{1+\alpha}\right) \right].
\]

Expanding the last logarithm,

\[
\log\left(\frac{1-\alpha}{1+\alpha}\right) = \log(1-\alpha)-\log(1+\alpha),
\]

all terms involving $\log(1+\alpha)$ cancel, leaving

\[
g(\alpha) = -1-\frac{1-\alpha}{\alpha}\log(1-\alpha), \qquad 0<\alpha<1. \tag{7}
\]

{\color{blue}
\paragraph*{Alternative derivation for $0<\alpha<1$}

The quadratic transformation in
\href{\detokenize{https://dlmf.nist.gov/15.8.E15}}{DLMF Section 15.8.15} is

\[
{}_2F_1(a,b;a-b+1;x)
=(1+x)^{-a}{}_2F_1\left(\frac a2,\frac{a+1}{2};a-b+1;
\frac{4x}{(1+x)^2}\right),
\qquad |x|<1.
\]

Taking $a=1-n$, $b=-n$, and $x=\alpha$ gives

\[
f(n,\alpha)={}_2F_1(1-n,-n;2;\alpha).
\]

By the Gauss series,

\[
f(n,\alpha)=\sum_{k=0}^{\infty}c_k(n)\alpha^k,
\qquad
c_k(n)=\frac{(1-n)_k(-n)_k}{(2)_k\,k!}.
\]

Termwise differentiation with respect to $n$ is justified as follows. For
fixed $\alpha\in(0,1)$ and $\delta>0$, the recurrence for $c_k$ gives

\[
\left|c_{k+1}(n)\alpha^{k+1}\right|
\leq \left|c_k(n)\alpha^k\right|
\alpha\frac{(k+1+\delta)(k+\delta)}{(k+2)(k+1)},
\qquad |n|\leq\delta,
\]

and the last factor tends uniformly to $\alpha<1$. The uniform ratio test
therefore gives locally uniform convergence in $n$. Since every $c_k$ is
entire, the Weierstrass theorem and Cauchy's integral formula imply that the
sum is holomorphic near $n=0$ and may be differentiated term by term there.

For $k\geq1$,

\[
(-n)_k=(-n)(1-n)\cdots(k-1-n),
\]

and hence

\[
\left.\frac{d}{dn}(-n)_k\right|_{n=0}=-(k-1)!,
\qquad
(1-n)_k\big|_{n=0}=k!.
\]

Since $(-n)_k|_{n=0}=0$, the derivative of $(1-n)_k$ does not contribute.
Using $(2)_k=(k+1)!$, we obtain

\[
c_0'(0)=0,
\qquad
c_k'(0)=-\frac{1}{k(k+1)},\qquad k\geq1.
\]

It follows that

\[
g(\alpha)=-\sum_{k=1}^{\infty}\frac{\alpha^k}{k(k+1)}.
\]

Finally, using $1/[k(k+1)]=1/k-1/(k+1)$ and
$-\log(1-\alpha)=\sum_{k=1}^{\infty}\alpha^k/k$, with both series
absolutely convergent for $0<\alpha<1$, gives

\[
\sum_{k=1}^{\infty}\frac{\alpha^k}{k(k+1)}
=1+\frac{1-\alpha}{\alpha}\log(1-\alpha).
\]

Therefore,

\[
g(\alpha)=-1-\frac{1-\alpha}{\alpha}\log(1-\alpha),
\qquad 0<\alpha<1,
\]

in agreement with (7).
}
\subparagraph*{5. Endpoint $\alpha=0$}

When $\alpha=0$, the hypergeometric argument is zero. Therefore,

\[
f(n,0) = 1^{n-1}\,{}_2F_1\left(\frac{1-n}{2},1-\frac n2;2;0\right) =1.
\]

Hence

\[
g(0)=0.
\]

This also agrees with the continuous limit because

\[
\log(1-\alpha)=-\alpha+O(\alpha^2),
\]

so

\[
-1-\frac{1-\alpha}{\alpha}\log(1-\alpha) =O(\alpha)\longrightarrow0.
\]

\subparagraph*{6. Endpoint $\alpha=1$}

At $\alpha=1$, the hypergeometric argument is $z=1$. Gauss's summation theorem, stated in \href{\detokenize{https://dlmf.nist.gov/15.4.E20}}{DLMF Section 15.4.20}, gives

\[
f(n,1) = 2^{n-1} \frac{ \Gamma\left(\frac12+n\right) }{ \Gamma\left(\frac{3+n}{2}\right) \Gamma\left(1+\frac n2\right) }.
\]

Since $f(0,1)=1$, logarithmic differentiation yields

\[
g(1) = \log2+\psi\left(\frac12\right) -\frac12\psi\left(\frac32\right) -\frac12\psi(1).
\]

The special values are

\[
\psi(1)=-\gamma, \qquad \psi\left(\frac12\right)=-\gamma-2\log2,
\]

as listed in \href{\detokenize{https://dlmf.nist.gov/5.4.E12}}{DLMF Section 5.4}, and the recurrence

\[
\psi(z+1)=\psi(z)+\frac1z
\]

from \href{\detokenize{https://dlmf.nist.gov/5.5.E2}}{DLMF Section 5.5.2} gives

\[
\psi\left(\frac32\right) = \psi\left(\frac12\right)+2.
\]

Consequently,

\[
g(1)=-1.
\]

This agrees with (7) under the standard limiting convention

\[
\lim_{x\to0^+}x\log x=0.
\]
